\documentclass[11pt]{article}
\usepackage[margin=0.68in]{geometry}
\usepackage{lmodern}
\usepackage{microtype}
\usepackage{enumitem}
\usepackage[table]{xcolor}
\usepackage{titlesec}
\setlength{\parindent}{0pt}
\setlength{\parskip}{0.25em}
\setlist[itemize]{leftmargin=1.3em,itemsep=0.06em,topsep=0.1em}
\titleformat{\section}{\large\bfseries}{}{0pt}{}
\titlespacing*{\section}{0pt}{0.55em}{0.2em}
\pagestyle{plain}
\newenvironment{playpassage}{\begin{quote}\small\setlength{\parskip}{0.05em}\setlength{\parindent}{0pt}}{\end{quote}}
\newcommand{\speaker}[1]{\textbf{#1:}}
\begin{document}
\begin{center}
{\Large\bfseries U1L10SP --- Text Supplement}\par
{\LARGE\bfseries Danger, Necessity, and Evidence}\par
{\large Two scenes for close reading; the whole play remains available}
\end{center}
\fcolorbox{black}{gray!8}{\begin{minipage}{0.94\linewidth}\small
\textbf{How to use this supplement}

You may use \textbf{anything accurate from \textit{Julius Caesar}} in your paper. The passages below are lines worth focusing on, not limits on your evidence. Quote exact words when they are printed here. You may accurately paraphrase other events you remember from watching the play. Do not invent quotations from memory.
\end{minipage}}

\section*{Passage 1: Act 2, Scene 1 --- Brutus in the Orchard}
\textbf{Context:} Before joining the conspiracy, Brutus reasons alone about Caesar. Caesar has not yet been crowned. Brutus considers what greater power might do to Caesar and what should be done now.
\begin{playpassage}
\speaker{Brutus} It must be by his death: and for my part,\par
I know no personal cause to spurn at him,\par
But for the general. He would be crown'd:\par
How that might change his nature, there's the question.\par
It is the bright day that brings forth the adder;\par
And that craves wary walking. Crown him? That;\par
And then, I grant, we put a sting in him,\par
That at his will he may do danger with.\par
The abuse of greatness is, when it disjoins\par
Remorse from power: and, to speak truth of Caesar,\par
I have not known when his affections sway'd\par
More than his reason. But 'tis a common proof,\par
That lowliness is young ambition's ladder,\par
Whereto the climber-upward turns his face;\par
But when he once attains the upmost round,\par
He then unto the ladder turns his back,\par
Looks in the clouds, scorning the base degrees\par
By which he did ascend. So Caesar may.\par
Then, lest he may, prevent. And, since the quarrel\par
Will bear no colour for the thing he is,\par
Fashion it thus; that what he is, augmented,\par
Would run to these and these extremities:\par
And therefore think him as a serpent's egg\par
Which, hatch'd, would, as his kind, grow mischievous,\par
And kill him in the shell.
\end{playpassage}
\section*{Words in context}
\begin{itemize}
\item \textbf{spurn at}: reject or act against; \textbf{for the general}: for the public good;
\item \textbf{remorse}: pity or moral restraint; \textbf{affections}: passions or desires;
\item \textbf{common proof}: common experience; \textbf{base degrees}: lower steps;
\item \textbf{prevent}: act beforehand; \textbf{colour}: plausible justification;
\item \textbf{augmented}: increased in power; \textbf{mischievous}: harmful.
\end{itemize}

\newpage
\section*{Passage 2: Act 3, Scene 2 --- Antony Presses Caesar's Record}
\textbf{Context:} Caesar has been killed. Brutus has told Rome that Caesar was ambitious and that his death protected Roman freedom. Antony responds with examples from Caesar's public behavior. Decide for yourself how far these examples reach.
\begin{playpassage}
\speaker{Antony} He hath brought many captives home to Rome\par
Whose ransoms did the general coffers fill:\par
Did this in Caesar seem ambitious?\par
When that the poor have cried, Caesar hath wept:\par
Ambition should be made of sterner stuff:\par
Yet Brutus says he was ambitious;\par
And Brutus is an honourable man.\par
You all did see that on the Lupercal\par
I thrice presented him a kingly crown,\par
Which he did thrice refuse: was this ambition?
\end{playpassage}
\section*{Words in context}
\begin{itemize}
\item \textbf{ransoms}: payments made to free captives;
\item \textbf{general coffers}: Rome's public treasury;
\item \textbf{sterner stuff}: harder or less compassionate qualities;
\item \textbf{Lupercal}: a Roman festival at which Antony offered Caesar a crown;
\item \textbf{thrice}: three times.
\end{itemize}

\section*{The rest of the play is available to you}
You may also use accurate remembered events: Caesar's conduct before the assassination, the conspirators' stated motives, Brutus's funeral speech, Antony's persuasion of the crowd, the civil conflict that follows, or Brutus's later judgments. These are optional reminders, not conclusions.

Use exact quotation marks only for words printed in a text you have. Other events may be paraphrased accurately. In every case, explain whether the evidence shows present fact, future possibility, or the necessity of a particular action.
\end{document}
